\defaultCJKfontfeatures{}
% Prefer TrueType CJK fonts: xdvipdfmx embeds explicit ToUnicode maps for them.
% This avoids relying on a PDF viewer's external Adobe CJK character maps.
\IfFontExistsTF{SimSun}{%
  \setCJKmainfont[BoldFont=SimHei,ItalicFont=KaiTi,BoldItalicFont=KaiTi]{SimSun}
  \setCJKsansfont[BoldFont=Microsoft YaHei Bold]{Microsoft YaHei}
  \setCJKmonofont{SimSun}
}{%
  \setCJKmainfont[BoldFont=FandolSong-Bold,ItalicFont=FandolKai-Regular]{FandolSong-Regular}
  \setCJKsansfont[BoldFont=FandolHei-Bold]{FandolHei-Regular}
  \setCJKmonofont{FandolFang-Regular}
}
\usepackage[a4paper,inner=24mm,outer=21mm,top=24mm,bottom=24mm,headheight=16pt]{geometry}
\usepackage{lmodern}
\usepackage{amsmath,amssymb,mathtools,bm,mathrsfs}
\usepackage[dvipsnames]{xcolor}
\usepackage[most]{tcolorbox}
\usepackage{tikz,pgfplots,booktabs,array,longtable,tabularx,enumitem,fancyhdr,graphicx}
\usepackage{titletoc,multicol,needspace}
\usepackage{fvextra}
\usepackage{hyperref,cleveref}
\fvset{fontsize=\footnotesize,breaklines=true,breakanywhere=false,numbers=left,
  breaksymbolleft=\textcolor{Rust}{$\hookrightarrow$},
  numbersep=5pt,xleftmargin=5mm,baselinestretch=1.05,
  frame=leftline,rulecolor=\color{Teal!45}}
\pgfplotsset{compat=1.18}
\definecolor{Ink}{HTML}{19364A}
\definecolor{Teal}{HTML}{227C81}
\definecolor{Gold}{HTML}{A87830}
\definecolor{Rust}{HTML}{9B483A}
\definecolor{CoverTeal}{HTML}{20ADAA}
\definecolor{CoverGold}{HTML}{DDBD35}
\hypersetup{colorlinks=true,linkcolor=Teal,urlcolor=Teal,citecolor=Teal,
  bookmarksnumbered=true,pdftitle={Python 网络爬虫：从请求到数据},
  pdfauthor={雪落长安无声},pdfsubject={Python 网络爬虫学习笔记}}
\ctexset{chapter={format=\Huge\bfseries\color{Ink},beforeskip=10pt,afterskip=22pt},section={format=\Large\bfseries\color{Teal}},subsection={format=\large\bfseries\color{Ink}}}
\setcounter{tocdepth}{2}
\setcounter{secnumdepth}{2}
% openany also applies when switching from front matter to main matter.
\let\cleardoublepage\clearpage
\pagestyle{fancy}\fancyhf{}
\fancyhead[LE,RO]{\small\color{Ink}\thepage}
\fancyhead[LO]{\small\nouppercase{\leftmark}}
\fancyhead[RE]{\small Python 网络爬虫：从请求到数据}
\renewcommand{\headrulewidth}{0.4pt}
\fancypagestyle{plain}{\fancyhf{}\fancyfoot[C]{\small\thepage}\renewcommand{\headrulewidth}{0pt}}
\tcbset{enhanced,breakable,boxrule=0.45pt,arc=1mm,left=2mm,right=2mm,top=1.2mm,bottom=1.2mm,before skip=8pt,after skip=8pt,fonttitle=\bfseries,colback=white,colframe=Teal}
\newcounter{statement}[chapter]\renewcommand{\thestatement}{\thechapter.\arabic{statement}}
\newenvironment{definition}[1]{\refstepcounter{statement}\begin{tcolorbox}[title={定义\thestatement\quad #1},colback=Teal!3]}{\end{tcolorbox}}
\newenvironment{theorem}[1]{\refstepcounter{statement}\begin{tcolorbox}[title={定理\thestatement\quad #1},colback=Ink!3,colframe=Ink]}{\end{tcolorbox}}
\newenvironment{lemma}[1]{\refstepcounter{statement}\begin{tcolorbox}[title={引理\thestatement\quad #1}]}{\end{tcolorbox}}
\newenvironment{corollary}[1]{\refstepcounter{statement}\begin{tcolorbox}[title={推论\thestatement\quad #1}]}{\end{tcolorbox}}
\newenvironment{proof}{\begin{tcolorbox}[title=原理与推导,colframe=Ink,colbacktitle=Ink,coltitle=white,colback=Ink!2]}{\end{tcolorbox}}
\newenvironment{method}[1]{\begin{tcolorbox}[title={方法卡\quad #1},colframe=Teal,colback=Teal!3]}{\end{tcolorbox}}
\newenvironment{warning}{\begin{tcolorbox}[title=常见错误与边界,colframe=Rust,colback=Rust!3]}{\end{tcolorbox}}
\newenvironment{examtip}{\begin{tcolorbox}[title=实践提示,colframe=Gold,colback=Gold!4]}{\end{tcolorbox}}
\newenvironment{extension}{\begin{tcolorbox}[title=进阶实践,colframe=Gold,colback=Gold!4]}{\end{tcolorbox}}
\newcounter{example}[chapter]\renewcommand{\theexample}{\thechapter.\arabic{example}}
\newenvironment{example}[2]{\refstepcounter{example}\begin{tcolorbox}[title={例\theexample\quad #1\quad #2},title after break={例\theexample\quad #1（续）},colframe=Ink]}{\end{tcolorbox}}
\newcounter{exercise}[chapter]\renewcommand{\theexercise}{\thechapter.\arabic{exercise}}
\newcommand{\dd}{\mathop{}\!\mathrm{d}}
\newcommand{\e}{\mathrm{e}}
\newcommand{\R}{\mathbb{R}}
\newcommand{\grad}{\nabla}
\DeclareMathOperator{\diver}{div}
\DeclareMathOperator{\rot}{rot}
\DeclareMathOperator{\sgn}{sgn}
\newcommand{\abs}[1]{\left|#1\right|}
\setlength{\parindent}{2em}\setlength{\parskip}{3pt}
\setlist{nosep,topsep=4pt,leftmargin=2em}
\renewcommand{\arraystretch}{1.25}
\emergencystretch=2em
\raggedbottom

% Printed contents follows the compact two-column chapter/section layout.
% Subsections remain available through the PDF bookmarks.
\definecolor{ContentsBlue}{HTML}{3578C4}
\hypersetup{bookmarksdepth=2}
\ctexset{chapter/tocline={\CTEXifname{\protect\numberline{\thechapter}}{}#2}}
\titlecontents{chapter}[1.7em]
  {\Needspace{3\baselineskip}\addvspace{6pt}\bfseries\color{black}}
  {\contentslabel{1.7em}}{}
  {\hfill\contentspage}[\addvspace{1pt}]
\titlecontents{section}[3.4em]
  {\color{black}}
  {\contentslabel{2.6em}}{}
  {\hspace{.5em}\titlerule*[.7pc]{.}\contentspage}
\contentsmargin{2.5em}
\makeatletter
\newcommand{\booktableofcontents}{%
  \clearpage
  \begingroup
  \pagestyle{plain}\markboth{目录}{目录}%
  \hypersetup{linkcolor=black}%
  \phantomsection\pdfbookmark[0]{目录}{book.contents}%
  \setlength{\parskip}{0pt}%
  \setlength{\columnsep}{8mm}\setlength{\columnseprule}{0pt}%
  \setlength{\multicolsep}{0pt}%
  \vspace*{-2mm}%
  \begin{center}{\fontsize{20}{26}\selectfont\bfseries\color{ContentsBlue}目录}\end{center}%
  \vspace{5mm}%
  \c@tocdepth=1\relax
  \begin{multicols}{2}%
    \fontsize{9.5}{14}\selectfont\raggedcolumns
    \@starttoc{toc}%
  \end{multicols}%
  \clearpage
  \endgroup
}
\makeatother
